\documentclass[12pt, a4paper]{article}

\usepackage[utf8]{inputenc}
\usepackage[T2A]{fontenc}
\usepackage[russian]{babel}

\usepackage{geometry}
\geometry{
    top=2cm,
    bottom=2cm,
    left=3cm,
    right=1.5cm
}

\usepackage{graphicx}
\usepackage{hyperref}
\usepackage{xcolor}
\usepackage{float}
\usepackage{booktabs}

\usepackage{listings}
\lstset{
    language=C++,
    basicstyle=\ttfamily\small,
    keywordstyle=\color{blue}\bfseries,
    commentstyle=\color{gray},
    stringstyle=\color{red},
    numbers=left,
    numberstyle=\tiny\color{gray},
    stepnumber=1,
    numbersep=5pt,
    frame=single,
    breaklines=true,
    captionpos=b,
    tabsize=4,
    showstringspaces=false,
    inputencoding=utf8,
    extendedchars=true,
    literate=
        {а}{{\cyra}}1 {б}{{\cyrb}}1 {в}{{\cyrv}}1 {г}{{\cyrg}}1
        {д}{{\cyrd}}1 {е}{{\cyre}}1 {ё}{{\"e}}1 {ж}{{\cyrzh}}1
        {з}{{\cyrz}}1 {и}{{\cyri}}1 {й}{{\cyrishrt}}1 {к}{{\cyrk}}1
        {л}{{\cyrl}}1 {м}{{\cyrm}}1 {н}{{\cyrn}}1 {о}{{\cyro}}1
        {п}{{\cyrp}}1 {р}{{\cyrr}}1 {с}{{\cyrs}}1 {т}{{\cyrt}}1
        {у}{{\cyru}}1 {ф}{{\cyrf}}1 {х}{{\cyrh}}1 {ц}{{\cyrc}}1
        {ч}{{\cyrch}}1 {ш}{{\cyrsh}}1 {щ}{{\cyrshch}}1 {ъ}{{\cyrhrdsn}}1
        {ы}{{\cyrery}}1 {ь}{{\cyrsoft}}1 {э}{{\cyre}}1 {ю}{{\cyryu}}1
        {я}{{\cyrya}}1
}

\usepackage{setspace}
\onehalfspacing

\begin{document}

\begin{titlepage}
    \begin{center}
        \Large\textbf{МИНИСТЕРСТВО НАУКИ И ВЫСШЕГО ОБРАЗОВАНИЯ РОССИЙСКОЙ ФЕДЕРАЦИИ} \\
        \vspace{0.3cm}
        \large Федеральное государственное автономное образовательное учреждение высшего образования \\
        \textbf{«Московский авиационный институт (национальный исследовательский университет)»}
    \end{center}

    \vspace{3cm}

    \begin{center}
        \huge\textbf{ОТЧЁТ ПО ЛАБОРАТОРНОЙ РАБОТЕ № 3}
    \end{center}

    \vspace{1.5cm}

    \begin{center}
        \large\textbf{Дисциплина:} Операционные системы и архитектура компьютера \\
        \vspace{0.5cm}
        \large\textbf{Тема:} Управление потоками в ОС. Обеспечение синхронизации между потоками
    \end{center}

    \vspace{3cm}

    \begin{flushright}
        \begin{tabular}{ll}
            \textbf{Выполнил:} & студент группы М8О-203БВ-25 \\
                               & Воскресенский И. А. \\[1cm]
            \textbf{Проверил:} & Преподаватель \\
                               & Миронов Е. С. \\
        \end{tabular}
    \end{flushright}

    \vfill

    \begin{center}
        \large Москва, 2026 г.
    \end{center}
\end{titlepage}

\tableofcontents
\newpage

\section{Цель работы}

Приобретение практических навыков в управлении потоками в операционной системе и обеспечении синхронизации между ними. Разработка программы на языке C++ для вычисления детерминанта матрицы в многопоточном режиме с использованием пула потоков, ограниченного ключом запуска. Исследование зависимости ускорения и эффективности алгоритма от входных данных и количества потоков.

\newpage

\section{Краткое описание архитектуры решения}

Решение состоит из главного потока (main), который отвечает за парсинг аргументов командной строки, генерацию матрицы и управление жизненным циклом пула потоков. 

Пул потоков (\texttt{ThreadPool}) содержит фиксированное количество рабочих потоков, задаваемое ключом \texttt{-t}. Потоки синхронизируются с помощью мьютекса (\texttt{std::mutex}) для защиты общей очереди задач и условной переменной (\texttt{std::condition\_variable}) для эффективного ожидания новых задач без активного потребления CPU.

Вычисление детерминанта выполняется методом разложения по строке (метод Лапласа). Для предотвращения исчерпания памяти из-за факториальной сложности алгоритма введено ограничение глубины параллельной рекурсии (\texttt{MAX\_PARALLEL\_DEPTH}). Результаты подзадач собираются асинхронно через механизм \texttt{std::future}.

\subsection{Схема взаимодействия компонентов}

Поток данных и управляющих сообщений:
\begin{enumerate}
    \item \texttt{Main} $\xrightarrow{\text{создает}}$ \texttt{ThreadPool(N)}
    \item \texttt{Main} $\xrightarrow{\text{отправляет задачу}}$ \texttt{Очередь задач (Mutex + CondVar)}
    \item \texttt{Очередь} $\xrightarrow{\text{будит}}$ \texttt{Worker Thread}
    \item \texttt{Worker Thread} $\xrightarrow{\text{возвращает результат}}$ \texttt{std::future} $\xrightarrow{\text{сборка}}$ \texttt{Main}
\end{enumerate}

\begin{figure}[H]
    \centering
    \fbox{
        \begin{minipage}{0.8\textwidth}
            \centering
            \texttt{Main Thread} $\rightarrow$ \texttt{ThreadPool (Queue)} $\rightarrow$ \texttt{Worker Threads (1..N)} \\
            \vspace{0.5cm}
            \texttt{Worker Threads} $\rightarrow$ \texttt{std::future} $\rightarrow$ \texttt{Main Thread}
        \end{minipage}
    }
    \caption{Архитектура решения}
    \label{fig:arch}
\end{figure}

\newpage

\section{Листинг ключевых фрагментов кода}

Полные исходные файлы доступны в репозитории: \texttt{https://github.com/fricexxx/osi-lab3}

\subsection{Инициализация пула потоков и синхронизация}

Фрагмент демонстрирует создание рабочих потоков и механизм их ожидания с использованием мьютекса и условной переменной. Это обеспечивает потокобезопасный доступ к очереди задач.

\begin{lstlisting}[caption={Конструктор ThreadPool и механизм ожидания}, label={lst:pool}]
ThreadPool(size_t numThreads) : stop(false) {
    for (size_t i = 0; i < numThreads; ++i) {
        workers.emplace_back([this] {
            while (true) {
                std::function<void()> task;
                {
                    std::unique_lock<std::mutex> lock(this->queueMutex);
                    this->condition.wait(lock, [this] { 
                        return this->stop || !this->tasks.empty(); 
                    });
                    if (this->stop && this->tasks.empty()) return;
                    task = std::move(this->tasks.front());
                    this->tasks.pop();
                }
                task();
            }
        });
    }
}
\end{lstlisting}

\newpage

\subsection{Параллельное вычисление детерминанта}

Фрагмент показывает отправку задач в пул и ограничение глубины рекурсии. Использование \texttt{std::shared\_ptr} для миноров предотвращает избыточное копирование данных при захвате лямбдой.

\begin{lstlisting}[caption={Рекурсивное вычисление с ограничением глубины}, label={lst:det}]
double Matrix::determinant(ThreadPool& pool, size_t maxThreads, size_t depth) const {
    if (maxThreads <= 1 || n < 4 || depth >= 3) {
        return determinantSequential();
    }
    
    std::vector<std::future<double>> futures;
    for (size_t j = 0; j < n; ++j) {
        double cofactor = (j % 2 == 0 ? 1.0 : -1.0) * data[j];
        auto minorPtr = std::make_shared<Matrix>(getMinor(0, j));
        
        futures.push_back(pool.enqueue(
            [cofactor, minorPtr, &pool, maxThreads, depth]() {
                return cofactor * minorPtr->determinant(pool, maxThreads, depth + 1);
            }
        ));
    }
    
    double det = 0.0;
    for (auto& f : futures) det += f.get();
    return det;
}
\end{lstlisting}

\newpage

\section{Результаты работы}

\subsection{Консольный вывод и исследование эффективности}

Ниже приведен вывод программы при запуске с матрицей $9 \times 9$ и разным количеством потоков. Различные TID подтверждают конкурентное выполнение.

\begin{verbatim}
$ ./determinant -t 4 -s 9
[Main PID=6323, TID=6323] Starting. Max threads: 4, Matrix size: 9
[Main PID=6323, TID=6323] Hardware concurrency: 8
[Main PID=6323, TID=6323] Determinant = -1.4532e+07
[Main PID=6323, TID=6323] Computation time: 142 ms
[Main PID=6323, TID=6323] Thread pool size: 4
\end{verbatim}

\textbf{Таблица 1. Зависимость времени, ускорения и эффективности от числа потоков (матрица 9x9)}

\begin{table}[H]
    \centering
    \begin{tabular}{@{}cccc@{}}
        \toprule
        \textbf{Кол-во потоков (N)} & \textbf{Время (мс)} & \textbf{Ускорение (S)} & \textbf{Эффективность (E)} \\ \midrule
        1 & 450 & 1.00 & 1.00 \\
        2 & 240 & 1.87 & 0.93 \\
        4 & 142 & 3.16 & 0.79 \\
        8 & 115 & 3.91 & 0.48 \\ \bottomrule
    \end{tabular}
    \caption{Результаты бенчмарка}
    \label{tab:benchmark}
\end{table}

\newpage

\section{Выводы}

В ходе выполнения лабораторной работы была реализована многопоточная программа на языке C++ для вычисления детерминанта матрицы с использованием паттерна «Пул потоков». Синхронизация потоков обеспечена через мьютексы и условные переменные, что позволило избежать состояний гонки и избыточного потребления CPU. 

Экспериментальное исследование показало, что параллельное выполнение дает значительное ускорение, однако эффективность падает при чрезмерном увеличении числа потоков из-за накладных расходов на синхронизацию и факториальной сложности алгоритма. Освоены практические навыки управления потоками ОС, асинхронного программирования и профилирования многопоточных приложений.

\end{document}